\subsection*{Changing the endpoint without changing the response sources}

Tree ensembles are effective interpolators, but their predictions are averages of
training-leaf values. They cannot continue a trend beyond their output range simply
because a query contains more atoms. We therefore tested whether the same frozen
response descriptors could support endpoints with different ways of representing
and combining molecular information (Table~\ref{tab:endpoints}). This comparison
includes residual multilayer perceptrons (MLPs), separate structure and response
branches, a linear predictor with a neural residual, message-passing networks with
size-sensitive readouts, frozen MoLFormer representations with neural heads, TabM
parameter-sharing ensembles, and TabPFN-3.5 in-context prediction. TabM shares most
weights across multiple neural predictions; TabPFN uses a pretrained transformer
conditioned on labelled training rows rather than fitting a new network by gradient
descent.\cite{tabm2025,tabpfn2025,tabpfn35}
Supplementary Figs.~8--10 explain the distinct mechanisms. The MoLFormer--ExtraTrees
control changes the representation without removing the tree endpoint.

The comparison is retrospective: these evaluation sets were already known during
method development. Within each outer partition, configuration and training duration
were chosen using only its training-side validation data. We completed longer-training,
representation and width challenges, locked the selection policy, and then fitted
the final models. The policy selected TabPFN in all five ARROW folds and the external
fit. The selected input recipe can differ by partition; the external TabPFN fit
uses RDKit descriptors, a frozen MoLFormer encoding and the responses.
The no-response controls inherit each selected recipe; they are matched ablations,
not separately optimized structural baselines. Original shuffled-response and
chemical-separation results remain evidence about ExtraTrees, not about these new heads.

\begin{table}[t]
\centering\small
\caption{\textbf{Endpoint mechanisms and their measured accuracy.}
All entries include the 15 response descriptors. MAE is in kcal mol$^{-1}$.
The first three columns use the original evaluation identities; Strict-97 is nested
within External-220. Neural and TabPFN predictions average five seeds, tree
controls three. The final column is a separate exploratory size-held-out refit on
1,160 training and 205 held-out molecules, using three seeds for every method.
A dash means not evaluated, not zero error. Supplementary Table~13 includes
matched no-response controls and common-three-seed results.}
\label{tab:endpoints}
\begin{tabular}{@{}p{3.65cm}p{4.7cm}rrrr@{}}
\toprule
Endpoint & Main mechanism & ARROW & External & Strict & Size \\
\midrule
ExtraTrees & Leaf averages & 0.202 & 1.153 & 1.536 & 1.772 \\
Residual MLP & Residual connections & 0.340 & 1.428 & 2.016 & 1.813 \\
Dual branch & Separate input branches & 0.288 & 1.234 & 1.616 & 1.560 \\
Ridge + neural residual & Linear trend + nonlinear residual & 0.336 & 1.172 & 1.429 & 1.391 \\
Message passing & Size-sensitive graph readout & 0.288 & 1.168 & 1.505 & 1.729 \\
MoLFormer + neural head & Frozen sequence representation & 0.273 & 1.111 & 1.450 & 1.609 \\
TabM & Shared-parameter ensemble & 0.291 & 1.132 & 1.454 & 1.389 \\
TabPFN-3.5 & Pretrained context prediction & 0.199 & 1.134 & 1.493 & 1.318 \\
MoLFormer + ExtraTrees & Representation-only change & 0.234 & 1.174 & 1.564 & -- \\
\bottomrule\end{tabular}

\end{table}

TabPFN reaches MAEs of 0.199, 1.134 and 1.493 kcal mol$^{-1}$ on ARROW-85,
External-220 and Strict-97, respectively. The corresponding ExtraTrees errors are
0.202, 1.153 and 1.536. The ARROW paired difference is $-0.003$ kcal mol$^{-1}$
(95\% interval, $-0.054$ to $0.066$); all three comparison intervals include zero.
The point estimates are close, but do not establish statistical superiority or
equivalence. Nor is response augmentation uniformly beneficial within TabPFN:
its no-response errors are 0.218, 1.131 and 1.347, and all three paired
response-minus-no-response intervals include zero (Fig.~\ref{fig:endpoints}).

Other architectures show different strengths. With a neural head, response
augmentation lowers MoLFormer errors from 0.465 to 0.273 on ARROW, 1.332 to 1.111
externally, and 1.685 to 1.450 on Strict-97; each paired interval is below zero.
Graph and TabM endpoints also show response benefits on both external evaluations,
although their ARROW errors remain about 0.29. Simply substituting a residual MLP
does not retain the original accuracy: its ARROW error is 0.340. A width-adaptive
selection variant reaches 1.090 externally, but its difference from ExtraTrees is
unresolved ($-0.062$; 95\% interval, $-0.157$ to $0.027$). This variant is a
training-side selection rule that can choose different configurations across
partitions, not a single additional architecture. Complete variants, unsuccessful
choices and paired intervals appear in Supplementary Tables~13--14 and Data~6.

\begin{figure}[p]
\centering
\includegraphics[width=\textwidth]{F7_endpoint_comparison.pdf}
\caption{\textbf{The response contribution depends on the endpoint architecture.}
\textbf{a--c}, Paired MAE changes from adding the same 15 response descriptors to
each selected endpoint recipe on ARROW-85, External-220 and Strict-97.
Negative values favor response augmentation. Bars are descriptive 95\% percentile
intervals from 100,000 paired molecule resamples, conditional on the fitted models
and without multiplicity correction. They are not uncertainty intervals for
individual molecular predictions. Strict-97 is nested within External-220.
The no-response controls inherit the full model's recipe. Cross-architecture
accuracy and seed counts are given in Table~\ref{tab:endpoints}; family-specific
configuration choices are retained in Supplementary Data~6.}
\label{fig:endpoints}
\end{figure}

\subsection*{Size transfer requires a separate test}

To test accuracy beyond the endpoint's training-size range, we held out the largest 15\% of
molecules within each endpoint source and refitted every evaluated model, including
ExtraTrees, on the remaining 1,160 labels. The 205 held-out molecules contain
8--37 heavy atoms; 188 exceed the largest training molecule, which has 17.
The overlapping lower end results from applying the rule within each source.
These labels were available during earlier random-fold development, so this is
an exploratory size diagnostic, not an untouched prospective test. The response
surrogates remain fixed: this comparison tests endpoint transfer conditional on
their predictions, not size extrapolation with all response-source exposure removed.

TabPFN reduces the size-held-out MAE from 1.772 for ExtraTrees to 1.318
kcal mol$^{-1}$ (paired difference $-0.453$; 95\% interval, $-0.648$ to $-0.281$;
Fig.~\ref{fig:size}a). Here its response descriptors also help: without them the
error is 1.959 (response-minus-control difference $-0.640$; interval, $-0.847$ to
$-0.451$). TabM, the linear-plus-neural model and the dual-branch model also reduce
error relative to ExtraTrees, with paired intervals below zero. These results
identify useful alternatives for this size split; they do not imply the same ordering
for every new chemical family (Supplementary Table~15).

Capped glycine and alanine series ask a complementary question about prediction
shape. ExtraTrees approaches plateaus near $-23$ and $-25$ kcal mol$^{-1}$, whereas
several neural endpoints and TabPFN continue to change over the tested range
(Fig.~\ref{fig:size}b,c). For Ac-(Gly)$_n$-NHMe, TabPFN predicts $-40.802$ at
$n=6$ and $-67.194$ at $n=12$, compared with $-23.005$ and $-23.440$ for
ExtraTrees. No experimental peptide values are available here: these are predictions,
not errors, and a more negative or more linear curve is not evidence of greater
physical accuracy. The finite outer bound on the mean prediction is approximately
$[-32.595,5.757]$ for the released trees and $[-658.495,649.922]$ for the fitted
TabPFN ensemble. The latter is much wider but still finite; neither architecture
guarantees asymptotic extensivity.

\begin{figure}[p]
\centering
\includegraphics[width=\textwidth]{F8_size_and_stress.pdf}
\caption{\textbf{Labelled size transfer and unlabelled series answer different questions.}
\textbf{a}, MAE on the 205-molecule size holdout after refitting on the same 1,160
training labels. Filled and open marks denote models with and without responses.
All predictions average three seeds. This diagnostic reuses development labels
under a different split. \textbf{b,c}, Predictions along capped glycine and alanine
series, using the full-data models with responses. Neural and TabPFN curves
average five seeds; ExtraTrees averages three. These curves have no peptide
reference values and therefore do not rank physical accuracy. The abscissa is the
number of residues between the acetyl and N-methylamide caps; $n=0$ is
N-methylacetamide. Colors identify the same families in both panels.}
\label{fig:size}
\end{figure}

Linear alkanes expose an additional limitation. TabPFN remains non-monotonic between
C10 and C14 (3.348 and 3.046 kcal mol$^{-1}$), despite its changed peptide behavior.
Moreover, 27 of the 30 tested alkanes occur in endpoint training, so this series is
not an independent extrapolation benchmark. A source audit confirmed that the
low C14 and C16 training values were copied faithfully from their source; removing
these two labels raises the corresponding tree predictions. We retain the original
labels and report the sensitivity rather than substitute an assumed linear trend
(Supplementary Methods~S20).

Finally, the native TabPFN predictive distributions do not yet provide calibrated
reliability. Equal mixtures of the five fitted distributions give nominal 80\%
and 95\% intervals with only 58.2\% and 80.0\% empirical coverage on External-220;
the corresponding Strict-97 coverages are 55.7\% and 76.3\%. These are predictive
intervals, not the spread across seeds. No calibration or model reselection was
performed on these outcomes (Supplementary Fig.~11 and Table~16).
